\chapter{Primitives et Intégrale}
\begin{objectifs}
	\begin{itemize}
		\item Calculer une primitive et une intégrale d’une fonction numérique
		\item Connaître quelques utilisations des intégrales de fonctions.
	\end{itemize}
\end{objectifs}

\begin{definition}
 $f$ est une fonction définie sur un intervalle $I$.\\
 Dire que $F$ est une \textbf{primitive} de $f$ signifie que $F$ est dérivable sur $I$ et \textbf{$F' = f$.}
\end{definition}

\begin{propriete}{}{}
	Toute fonction définie sur un interval $I$ admet des primitives sur $I$. 
\end{propriete}





% ============================================================
%  CHAPITRE : PRIMITIVES ET INTÉGRALE (suite)
%  Mazava Maths — Terminale S
%  Environnements : definition, propriete, methode, attention,
%                   astucebac, exemple, exercice, solution
% ============================================================

% ----------------------------------------------------------------
\section{Primitives des fonctions usuelles}
% ----------------------------------------------------------------

Le tableau suivant donne les primitives des fonctions usuelles
à connaître en Terminale~S.

\begin{center}
	\renewcommand{\arraystretch}{2.2}
	\begin{tabular}{|c|c|c|}
		\hline
		\textbf{Fonction} $f(x)$ & \textbf{Primitive} $F(x)$ & \textbf{Conditions} \\
		\hline
		$k$ \quad ($k \in \R$) & $kx$ & sur $\R$ \\
		\hline
		$x^n$ \quad ($n \in \N^*$) & $\dfrac{x^{n+1}}{n+1}$ & sur $\R$ \\
		\hline
		$x^n$ \quad ($n \in \Z,\ n \leq -2$) & $\dfrac{x^{n+1}}{n+1}$ & sur $]0;+\infty[$ ou $]-\infty;0[$ \\
		\hline
		$x^\alpha$ \quad ($\alpha \in \R,\ \alpha \neq -1$) & $\dfrac{x^{\alpha+1}}{\alpha+1}$ & sur $]0;+\infty[$ \\
		\hline
		$\dfrac{1}{x^2}$ & $-\dfrac{1}{x}$ & sur $]0;+\infty[$ ou $]-\infty;0[$ \\
		\hline
		$\dfrac{1}{\sqrt{x}}$ & $2\sqrt{x}$ & sur $]0;+\infty[$ \\
		\hline
		$\sqrt{x}$ & $\dfrac{2}{3}x\sqrt{x}$ & sur $]0;+\infty[$ \\
		\hline
		$\cos x$ & $\sin x$ & sur $\R$ \\
		\hline
		$\sin x$ & $-\cos x$ & sur $\R$ \\
		\hline
		$\dfrac{1}{\cos^2 x}$ & $\tan x$ & sur $\left]-\dfrac{\pi}{2};\dfrac{\pi}{2}\right[$ \\
		\hline
	\end{tabular}
\end{center}

\begin{attention}
	Une primitive n'est \textbf{jamais unique} : si $F$ est une primitive de $f$,
	alors $F + C$ (avec $C \in \R$) est aussi une primitive de $f$.
	On dit que les primitives de $f$ forment une \textbf{famille}.
\end{attention}

% ----------------------------------------------------------------
\section{Opérations sur les primitives}
% ----------------------------------------------------------------

\begin{propriete}{Linéarité}{}
	Soient $f$ et $g$ deux fonctions admettant des primitives $F$ et $G$
	sur un intervalle $I$, et $\lambda, \mu \in \R$. Alors :
	\begin{itemize}
		\item $F + G$ est une primitive de $f + g$ sur $I$.
		\item $\lambda F$ est une primitive de $\lambda f$ sur $I$.
		\item $\lambda F + \mu G$ est une primitive de $\lambda f + \mu g$ sur $I$.
	\end{itemize}
\end{propriete}

\begin{propriete}{Primitive d'une fonction composée — cas usuels}{}
	Soient $u$ une fonction dérivable sur $I$ et $n \in \Z \setminus \{-1\}$,
	$\alpha \in \R \setminus \{-1\}$. On a :
	
	\begin{center}
		\renewcommand{\arraystretch}{2.2}
		\begin{tabular}{|c|c|}
			\hline
			\textbf{Fonction} $f$ & \textbf{Primitive} $F$ \\
			\hline
			$u' \cdot u^n$ & $\dfrac{u^{n+1}}{n+1}$ \\
			\hline
			$u' \cdot u^\alpha$ & $\dfrac{u^{\alpha+1}}{\alpha+1}$ \\
			\hline
			$\dfrac{u'}{u^2}$ & $-\dfrac{1}{u}$ \\
			\hline
			$\dfrac{u'}{\sqrt{u}}$ & $2\sqrt{u}$ \\
			\hline
			$u' \cos u$ & $\sin u$ \\
			\hline
			$u' \sin u$ & $-\cos u$ \\
			\hline
		\end{tabular}
	\end{center}
\end{propriete}

\begin{methode}
	Pour reconnaître une primitive de la forme $u' \cdot u^n$ :
	\begin{enumerate}
		\item Identifier la fonction $u(x)$ à l'intérieur.
		\item Calculer $u'(x)$.
		\item Vérifier que la fonction à intégrer est bien du type $u' \cdot u^n$
		(au coefficient multiplicatif près).
		\item Appliquer la formule, puis ajuster le coefficient.
	\end{enumerate}
\end{methode}

\begin{exemple}
	Trouver une primitive de $f(x) = 3x^2(x^3+1)^4$.
	
	On pose $u(x) = x^3 + 1$, donc $u'(x) = 3x^2$.
	
	On reconnaît $f(x) = u'(x) \cdot [u(x)]^4$, donc :
	\[
	F(x) = \frac{[u(x)]^5}{5} = \frac{(x^3+1)^5}{5}.
	\]
\end{exemple}

\begin{exemple}
	Trouver une primitive de $g(x) = \dfrac{2x}{\sqrt{x^2+3}}$.
	
	On pose $u(x) = x^2 + 3$, donc $u'(x) = 2x$.
	
	On reconnaît $g(x) = \dfrac{u'(x)}{\sqrt{u(x)}}$, donc :
	\[
	G(x) = 2\sqrt{u(x)} = 2\sqrt{x^2+3}.
	\]
\end{exemple}

% ----------------------------------------------------------------
\section{Primitive prenant une valeur donnée en un point}
% ----------------------------------------------------------------

\begin{propriete}{Unicité sous condition initiale}{}
	Soit $f$ une fonction continue sur un intervalle $I$,
	$x_0 \in I$ et $y_0 \in \R$.
	Il existe \textbf{une unique primitive} $F$ de $f$ sur $I$
	telle que $F(x_0) = y_0$.
\end{propriete}

\begin{methode}
	Pour trouver la primitive $F$ vérifiant $F(x_0) = y_0$ :
	\begin{enumerate}
		\item Calculer la primitive générale $F(x) = \ldots + C$.
		\item Écrire la condition : $F(x_0) = y_0$.
		\item Résoudre pour trouver $C$.
	\end{enumerate}
\end{methode}

\begin{exemple}
	Soit $f(x) = 3x^2 - 4x + 1$. Trouver la primitive $F$ de $f$
	telle que $F(1) = 2$.
	
	Primitive générale : $F(x) = x^3 - 2x^2 + x + C$.
	
	Condition : $F(1) = 1 - 2 + 1 + C = C = 2$, donc $C = 2$.
	
	\[
	\boxed{F(x) = x^3 - 2x^2 + x + 2.}
	\]
\end{exemple}

% ----------------------------------------------------------------
\section{Intégrale d'une fonction}
% ----------------------------------------------------------------

\begin{definition}
	Soit $f$ une fonction \textbf{continue} et \textbf{positive} sur $[a\,;b]$,
	et $F$ une primitive de $f$ sur $[a\,;b]$.
	On appelle \textbf{intégrale} de $f$ entre $a$ et $b$ le réel :
	\[
	\int_a^b f(x)\,\dd x = \Big[F(x)\Big]_a^b = F(b) - F(a).
	\]
	$a$ est la \textbf{borne inférieure}, $b$ la \textbf{borne supérieure}.
\end{definition}

\begin{remarque}
	Le résultat $F(b)-F(a)$ est \textbf{indépendant du choix} de la primitive $F$,
	car deux primitives ne diffèrent que d'une constante qui s'annule dans la soustraction.
\end{remarque}

\begin{propriete}{Propriétés de l'intégrale}{}
	Soient $f$ et $g$ continues sur $[a\,;b]$, et $\lambda \in \R$ :
	\begin{itemize}
		\item \textbf{Linéarité :}
		$\displaystyle\int_a^b \bigl(\lambda f(x)+g(x)\bigr)\,\dd x
		= \lambda\int_a^b f(x)\,\dd x + \int_a^b g(x)\,\dd x.$
		\item \textbf{Relation de Chasles :}
		$\displaystyle\int_a^b f(x)\,\dd x
		= \int_a^c f(x)\,\dd x + \int_c^b f(x)\,\dd x$
		pour tout $c \in [a\,;b]$.
		\item \textbf{Inversion des bornes :}
		$\displaystyle\int_a^b f(x)\,\dd x = -\int_b^a f(x)\,\dd x.$
		\item \textbf{Intégrale nulle :}
		$\displaystyle\int_a^a f(x)\,\dd x = 0.$
		\item \textbf{Positivité :}
		si $f(x) \geq 0$ sur $[a\,;b]$, alors
		$\displaystyle\int_a^b f(x)\,\dd x \geq 0.$
		\item \textbf{Comparaison :}
		si $f(x) \leq g(x)$ sur $[a\,;b]$, alors
		$\displaystyle\int_a^b f(x)\,\dd x \leq \int_a^b g(x)\,\dd x.$
	\end{itemize}
\end{propriete}

\begin{exemple}
	Calculer $\displaystyle I = \int_1^3 (2x^2 - 3x + 1)\,\dd x$.
	
	\[
	I = \left[\frac{2x^3}{3} - \frac{3x^2}{2} + x\right]_1^3
	= \left(\frac{54}{3} - \frac{27}{2} + 3\right)
	- \left(\frac{2}{3} - \frac{3}{2} + 1\right).
	\]
	\[
	= \left(18 - \frac{27}{2} + 3\right) - \left(\frac{2}{3} - \frac{3}{2} + 1\right)
	= \frac{21}{2} - \frac{1}{6}
	= \frac{63 - 1}{6} = \frac{62}{6} = \frac{31}{3}.
	\]
\end{exemple}

% ----------------------------------------------------------------
\section{Interprétation géométrique — Aire}
% ----------------------------------------------------------------

\begin{propriete}{Aire sous une courbe}{}
	Soit $f$ continue et \textbf{positive} sur $[a\,;b]$.
	L'aire (en unités d'aire) du domaine délimité par la courbe
	$\mathcal{C}_f$, l'axe des abscisses et les droites $x=a$, $x=b$ est :
	\[
	\mathcal{A} = \int_a^b f(x)\,\dd x.
	\]
\end{propriete}

% Figure TikZ — aire sous une courbe
\begin{center}
	\begin{tikzpicture}[>=Stealth, scale=1.0]
		\begin{scope}
			\clip (0.5,-0.1) rectangle (3.5,2.5);
			\fill[gray!25]
			(1,0) -- plot[domain=1:3, samples=60]
			(\x, {0.3*\x*\x - 0.2*\x + 0.5}) -- (3,0) -- cycle;
		\end{scope}
		\draw[->] (-0.2,0) -- (4.2,0) node[right] {$x$};
		\draw[->] (0,-0.2) -- (0,2.8) node[above] {$y$};
		\draw[thick, domain=0.3:3.8, samples=80]
		plot (\x, {0.3*\x*\x - 0.2*\x + 0.5});
		\draw[dashed] (1,0) node[below] {$a$} -- (1,{0.3-0.2+0.5});
		\draw[dashed] (3,0) node[below] {$b$} -- (3,{0.3*9-0.6+0.5});
		\node at (2, 0.6) {$\mathcal{A}$};
		\node[right] at (3.8, {0.3*3.8*3.8-0.2*3.8+0.5}) {$\mathcal{C}_f$};
	\end{tikzpicture}
\end{center}

\begin{propriete}{Aire entre deux courbes}{}
	Soient $f$ et $g$ continues sur $[a\,;b]$ avec $f(x) \geq g(x)$.
	L'aire du domaine compris entre $\mathcal{C}_f$ et $\mathcal{C}_g$ est :
	\[
	\mathcal{A} = \int_a^b \bigl[f(x) - g(x)\bigr]\,\dd x.
	\]
\end{propriete}

\begin{attention}
	Si $f$ change de signe sur $[a\,;b]$, il faut \textbf{découper} l'intervalle
	aux zéros de $f$ et sommer les valeurs absolues de chaque intégrale partielle.
\end{attention}

% ----------------------------------------------------------------
\section{Valeur moyenne}
% ----------------------------------------------------------------

\begin{definition}
	La \textbf{valeur moyenne} de $f$ continue sur $[a\,;b]$ est :
	\[
	\mu = \frac{1}{b-a}\int_a^b f(x)\,\dd x.
	\]
\end{definition}

% ----------------------------------------------------------------
\section{Intégration par parties (IPP)}
% ----------------------------------------------------------------

\begin{theoreme}[Intégration par parties]
	Soient $u$ et $v$ deux fonctions dérivables sur $[a\,;b]$,
	à dérivées continues. Alors :
	\[
	\boxed{\int_a^b u(x)\,v'(x)\,\dd x
		= \Bigl[u(x)\,v(x)\Bigr]_a^b - \int_a^b u'(x)\,v(x)\,\dd x.}
	\]
\end{theoreme}

\begin{methode}
	Pour appliquer une IPP :
	\begin{enumerate}
		\item Choisir $u$ et $v'$ de sorte que $u'$ soit \textbf{plus simple} que $u$
		et que $v$ soit \textbf{facile à calculer} depuis $v'$.
		\item Calculer $u'$ et $v$.
		\item Appliquer la formule.
		\item Calculer l'intégrale restante.
	\end{enumerate}
\end{methode}

\begin{astucebac}
	En pratique au BAC, l'IPP est souvent utilisée avec :
	$u = $ polynôme (donc $u'$ diminue le degré) et
	$v' = \cos x$, $\sin x$ ou $x^\alpha$ (donc $v$ reste simple).
\end{astucebac}

\begin{exemple}
	Calculer $\displaystyle I = \int_0^{\pi} x\cos x\,\dd x$.
	
	On pose $u(x) = x$ et $v'(x) = \cos x$,
	donc $u'(x) = 1$ et $v(x) = \sin x$.
	
	\[
	I = \Big[x\sin x\Big]_0^{\pi} - \int_0^{\pi} \sin x\,\dd x
	= 0 - \Big[-\cos x\Big]_0^{\pi}
	= -\bigl(-\cos\pi + \cos 0\bigr)
	= -\bigl(1 + 1\bigr) = -2.
	\]
\end{exemple}

% ----------------------------------------------------------------
\section{Volume d'un solide de révolution}
% ----------------------------------------------------------------

\begin{propriete}{Volume de révolution}{}
	Soit $f$ continue et positive sur $[a\,;b]$.
	Le volume du solide engendré par la rotation de $\mathcal{C}_f$
	autour de l'axe des abscisses est :
	\[
	V = \pi \int_a^b \bigl[f(x)\bigr]^2\,\dd x.
	\]
\end{propriete}

\begin{exemple}
	Calculer le volume du solide de révolution engendré par $f(x) = \sqrt{x}$
	sur $[0\,;4]$ autour de l'axe des abscisses.
	
	\[
	V = \pi\int_0^4 \bigl(\sqrt{x}\bigr)^2\,\dd x
	= \pi\int_0^4 x\,\dd x
	= \pi\left[\frac{x^2}{2}\right]_0^4
	= \pi \cdot \frac{16}{2}
	= 8\pi.
	\]
\end{exemple}

% ----------------------------------------------------------------
%  EXERCICES RÉSOLUS
% ----------------------------------------------------------------

\exercice{\etoilepleine\etoilevide\etoilevide}
Calculer les intégrales suivantes :
\[
\text{a)}\quad I_1 = \int_0^2 (3x^2+2x-1)\,\dd x
\qquad
\text{b)}\quad I_2 = \int_1^4 \frac{1}{\sqrt{x}}\,\dd x
\qquad
\text{c)}\quad I_3 = \int_0^\pi \sin x\,\dd x
\]

\solution

\textbf{a)}
\[
I_1 = \left[x^3 + x^2 - x\right]_0^2
= (8+4-2) - 0 = 10.
\]

\textbf{b)}
\[
I_2 = \left[2\sqrt{x}\right]_1^4
= 2\sqrt{4} - 2\sqrt{1} = 4 - 2 = 2.
\]

\textbf{c)}
\[
I_3 = \left[-\cos x\right]_0^\pi
= -\cos\pi + \cos 0 = 1 + 1 = 2.
\]

\bigskip
\exercice{\etoilepleine\etoilepleine\etoilevide}
\begin{enumerate}
	\item Trouver la primitive $F$ de $f(x) = 4x(x^2-1)^3$ telle que $F(0)=1$.
	\item Calculer l'aire du domaine délimité par la courbe $\mathcal{C}_f$
	où $f(x) = x^2 - x$, l'axe des abscisses et les droites $x=0$, $x=2$.
\end{enumerate}

\solution

\textbf{1.} On pose $u = x^2-1$, $u'=2x$.
On écrit $f(x) = 2 \cdot 2x \cdot (x^2-1)^3 = 2u'u^3$, donc :
\[
F(x) = 2 \cdot \frac{u^4}{4} + C = \frac{(x^2-1)^4}{2} + C.
\]
Condition $F(0)=1$ : $\dfrac{(-1)^4}{2}+C = 1 \Rightarrow C = \dfrac{1}{2}$.
\[
\boxed{F(x) = \frac{(x^2-1)^4}{2} + \frac{1}{2}.}
\]

\textbf{2.} On cherche les zéros de $f(x)=x^2-x=x(x-1)$ sur $[0;2]$ : $x=0$ et $x=1$.
\begin{itemize}
	\item Sur $[0;1]$ : $f(x) \leq 0$, donc l'aire partielle est
	$\displaystyle\mathcal{A}_1 = -\int_0^1(x^2-x)\,\dd x
	= -\left[\frac{x^3}{3}-\frac{x^2}{2}\right]_0^1
	= -\left(\frac{1}{3}-\frac{1}{2}\right) = \frac{1}{6}$.
	\item Sur $[1;2]$ : $f(x) \geq 0$, donc
	$\displaystyle\mathcal{A}_2 = \int_1^2(x^2-x)\,\dd x
	= \left[\frac{x^3}{3}-\frac{x^2}{2}\right]_1^2
	= \left(\frac{8}{3}-2\right)-\left(\frac{1}{3}-\frac{1}{2}\right)
	= \frac{5}{6}$.
\end{itemize}
\[
\mathcal{A} = \mathcal{A}_1 + \mathcal{A}_2 = \frac{1}{6}+\frac{5}{6} = 1 \text{ u.a.}
\]

\bigskip
\exercice{\etoilepleine\etoilepleine\etoilepleine}
Calculer $\displaystyle I = \int_0^{\frac{\pi}{2}} x\sin x\,\dd x$
par une intégration par parties, puis en déduire le volume du solide
de révolution engendré par $f(x)=\sin x$ sur $\left[0;\dfrac{\pi}{2}\right]$.

\solution

\textbf{IPP :} On pose $u=x$, $v'=\sin x$,
donc $u'=1$, $v=-\cos x$.
\[
I = \Big[-x\cos x\Big]_0^{\frac{\pi}{2}} + \int_0^{\frac{\pi}{2}}\cos x\,\dd x
= 0 + \Big[\sin x\Big]_0^{\frac{\pi}{2}}
= 1.
\]

\textbf{Volume :}
\[
V = \pi\int_0^{\frac{\pi}{2}}\sin^2 x\,\dd x
= \pi\int_0^{\frac{\pi}{2}}\frac{1-\cos 2x}{2}\,\dd x
= \frac{\pi}{2}\left[x - \frac{\sin 2x}{2}\right]_0^{\frac{\pi}{2}}
= \frac{\pi}{2}\cdot\frac{\pi}{2}
= \frac{\pi^2}{4}.
\]
